%% ma: L12-B10-0006
%% lop: 12
%% chuong: 3
%% bai: 10
%% dang: 12.10.2
%% loai: TN4
%% muc: TH
%% dap_an: "B"
%% nguon: "(NBV)-ĐỀ SỐ 14-MỖI NGÀY 1 ĐỀ THI 2025.pdf (D14, MỖI NGÀY 1 ĐỀ - Đề số 14), Phần I Câu 2"
%% kien_thuc: "Bài 10 (độ lệch chuẩn của mẫu số liệu ghép nhóm)"
%% trang_thai: chinh-thuc
%% ghi_chu: "Phương án A là phương sai mẫu hiệu chỉnh (chia cho n-1), phương án C là phương sai, phương án D là độ lệch chuẩn hiệu chỉnh. Cùng ý tưởng với dạng 12.10.2."
\begin{ex}
Thống kê chiều cao (đơn vị: cm) của các học sinh trong một lớp học ta có bảng số liệu sau:
\begin{center}
\small\setlength{\tabcolsep}{3pt}
\begin{tabular}{|c|c|c|c|c|c|c|}
\hline
Chiều cao & $[150;155)$ & $[155;160)$ & $[160;165)$ & $[165;170)$ & $[170;175)$ & $[175;180)$\\
\hline
Số học sinh & $1$ & $4$ & $10$ & $9$ & $4$ & $2$\\
\hline
\end{tabular}
\end{center}
Độ lệch chuẩn của mẫu số liệu trên bằng bao nhiêu? (Kết quả làm tròn đến hàng phần trăm)
\choice{$35{,}66$}{\True $5{,}87$}{$34{,}47$}{$5{,}97$}
\loigiai{Các giá trị đại diện: $152{,}5;\ 157{,}5;\ 162{,}5;\ 167{,}5;\ 172{,}5;\ 177{,}5$ với cỡ mẫu $n=30$.
Số trung bình $\bar x=\dfrac{1\cdot152{,}5+4\cdot157{,}5+10\cdot162{,}5+9\cdot167{,}5+4\cdot172{,}5+2\cdot177{,}5}{30}\approx165{,}33$.
Phương sai $s^2=\dfrac{\sum n_i(x_i-\bar x)^2}{30}\approx34{,}47$, độ lệch chuẩn $s=\sqrt{34{,}47}\approx5{,}87$.}
\end{ex}
